%% Insert your citation key below as it is in main.bib file to generate References in IEEE format.
\IEEEPARstart{T}{he} penetration of renewable energy sources (RES)
projected to increase over the coming decades, with gains highly powered by
solar and wind resources. The NREL 
has solar and wind energy increasing to $>$50\% share by 2050
\cite{lin_research_2020}. 
Solar photovoltaics (PV) and Type III and Type IV wind turbine generators
(WTG) are both inverter-based resources (IBRs), in contrast to
synchronous generators (SGs) of which the conventional power grid is
mostly composed.

Throughout the history of the power grid, reliability has been a prime concern.
A core aspect of power system reliability is frequency control, the ability
to maintain the AC voltage and current signals close to a given setpoint
(e.g. 60 Hz in the USA) \cite{anderson_power_2003}. If frequency deviates too far from its setpoint,
protection mechanisms of the grid, such as under-frequency load shedding
(UFLS) are triggered to prevent damage to equipment and to arrest further
instability \cite{denholm_inertia_2020}
Load shedding means that some customer(s) have lost power, a serious grid
reliability issue.


When the grid is near equilibrium, its frequency is closely related
to the real power transfer between its nodes. When a system fault, such as a
generator going offline, causes a power imbalance in the system, 
the rest of the generators increase power output to balance
generation and load. To transfer power between nodes requires the phase angle
between nodes in the system to shift, which results in transient frequency
deviations.

The rate at which frequency changes (RoCoF) is related to underlying physical
system from which the energy is derived. Conventionally,
thermal power plants use SGs to convert the underlying mechanical energy
contained in steam into electrical energy. The RoCoF is limited by
the mechanical {\bf inertia} of the SG as governed by the swing equation.

On the other hand, IBRs connect an energy resource such as a PV cell,
WTG, or a battery energy storage system (bESS) to the grid via a 
power electronic (PE) inverter which converts a DC source to an AC output.
The dynamics governing the energy resource as well as the control algorithm.
Because the IBR is not directly connected to any form of mechanical inertia,
IBR do not have inertia and that systems with high penetrations of IBRs
are {\it low-inertia} systems.

The control algorithms used by IBRs can be broadly divided into two strategies:
\begin{enumerate}
  \item Grid-following Inverters (GFL), and
  \item Grid-forming Inverters (GFM).
\end{enumerate}

While the IEEE has established standards for GFL converters \cite{photovoltaics_ieee_2018},
the definition of GFM remains debated within research communities.
To date, there is no IEEE standard on the definition of GFM.
NREL defines a GFM to be any inverter that does not operate with a phase-lock loop (PLL)
\cite{lin_research_2020}

There has been some debate on the security of IBRs. The authors
of \cite{milano_foundations_2018}
posit that "low inertia likely implies low security" since
lack of mechanical inertia means that the stability of the system
relies on the software implementation of control algorithms. Thus,
there is a growing body of research aimed at demonstrating the 
robustness of the control strategies deployed on IBRs.

Despite the informal definition, GFM control technology is
not new: since the 1990s, it has been researched extensively by
the microgrid research community \cite{lin_research_2020},
and implemented within smaller power systems. However, the 
integration of GFM systems into the power grid on a larger scale
has 

While the deployment of increasing penetrations of IBR to the large-scale
grid may not be gauranteed, it remains one of the currently most 
economically feasible paths to net-zero carbon emissions. Furthermore,
any potential deployment will necessarily be gradual process: the entirety of
SG-based generation cannot be replaced with generation based on IBRs
overnight. Therefore, it is of vital importance to study the 
dynamics of IBR-based frequency control in the variety of contexts
anticipated in the forthcoming energy transition. Since this
entails nothing less than a complete rearchitecting of the energy
grid, there are numerous open research questions in this field
which will be enumerated.

This literature review focuses on the challenges and oppurtunities along
the path to incorporating high levels of IBR penetration into the grid 
at a large scale. It is organized as follows: section \ref{sec:math-models}
summarizes the mathematical formulation for frequency control; section
\ref{sec:control-strategies} enumerates and classifies the many
of control schemes that have been proposed in literature; section
\ref{sec:commericial-demonstrations} summarizes the major commericial
projects which have been deployed using GFM technology to-date.


